\documentclass[12pt,a4paper]{article}
\usepackage[utf8]{inputenc}
\usepackage[T1]{fontenc}
\usepackage{mathptmx}
\usepackage{amsmath,amssymb}
\usepackage{graphicx}
\usepackage[margin=1in]{geometry}
\usepackage{booktabs}
\usepackage{hyperref}
\usepackage{natbib}
\usepackage{caption}
\usepackage{float}
\usepackage{setspace}
\usepackage{xcolor}
\usepackage{xurl}
\onehalfspacing
\hypersetup{colorlinks=true,linkcolor=blue!60!black,citecolor=blue!60!black,urlcolor=blue!60!black,breaklinks=true}

\title{\textbf{From Hypothesis to Reality: Scale-Dependent\\
Universality Classes in Earth's Seismotectonic Rhythms}\\[4pt]
\large A Random Matrix Theory Investigation, from Forward Modeling\\
to Empirical Validation Across Eight Global Fault Zones}

\author{Ruqing Chen\\
\textit{GUT Geoservice Inc., Montreal, Canada}\\
\texttt{ruqing@hotmail.com}}
\date{June 2026}

\begin{document}
\maketitle

\begin{abstract}
We test whether the temporal organization of earthquakes conforms to the universality classes of random matrix theory (RMT), following a complete scientific arc from hypothesis through forward modeling to empirical validation. \textbf{Part~I (hypothesis)}: motivated by elastic-rebound theory, we hypothesize that fault systems occupy a gradient of RMT classes---from Gaussian-Orthogonal/Unitary-Ensemble (GOE/GUE) level repulsion on simple locked faults to Poisson randomness in complex media---indexed by a repulsion parameter $\hat\beta$. Using forward physical modeling with synthetic catalogs calibrated to published Gutenberg--Richter parameters, we demonstrate the predicted gradient and characterize the estimator's bias and statistical power. \textbf{Part~II (empirical test)}: we then confront this hypothesis with real seismicity from eight globally distributed fault zones (USGS/ANSS, 1970--2024), processed through magnitude homogenization, completeness estimation, and dual declustering. The hypothesis is \emph{rejected at the mainshock scale}: declustered mainshock spacings collapse onto the Poisson distribution for all eight zones ($\langle r\rangle\approx0.39$), with no GOE signal. \textbf{Part~III (resolution)}: rather than a null result, this reveals a richer truth. The full catalogs are strongly clustered (coefficient of variation $1.1$--$3.8$; Gamma shape $\gamma=0.16$--$0.69$), and isolating the cluster process yields an Omori--Utsu decay exponent $p$ that does \emph{not} collapse and varies systematically with tectonic setting ($p\approx1.28$ for the dry Parkfield segment versus $0.6$--$0.9$ for fluid-rich subduction). We conclude that Earth's seismotectonic rhythms occupy three distinct universality classes at three temporal scales---Omori clustering (hours--days), Poisson noise (years--decades), and GOE repulsion at geological-epoch scales (Myr, from a companion study)---and that the Omori $p$-exponent, not a mainshock repulsion parameter, is the physically meaningful invariant distinguishing tectonic environments.

\medskip
\noindent\textbf{Keywords:} random matrix theory; earthquake statistics; Omori--Utsu law; declustering; scale-dependent universality; forward modeling; inter-event times; ETAS

\medskip
\noindent\textbf{Code/data:} \url{https://github.com/Ruqing1963/fault-rmt-phase-framework}\\
\textbf{Companion geological study:} \url{https://zenodo.org/records/20766310}
\end{abstract}

\section{Introduction}

Whether the timing of earthquakes carries ``memory'' is a central question in statistical seismology. Reid's elastic-rebound theory \citep{reid1910} predicts quasi-periodic recurrence on individual faults, which would manifest as a suppression of short inter-event times---the seismic analog of quantum-mechanical level repulsion described by random matrix theory (RMT) \citep{wigner1957,mehta2004}. The Dyson symmetry index $\beta$ quantifies this repulsion, with $\beta=0$ (Poisson, no memory), $\beta=1$ (Gaussian Orthogonal Ensemble, GOE), and $\beta=2$ (Gaussian Unitary Ensemble, GUE).

This paper presents a complete investigation of that question, structured as a scientific arc. We first develop, through forward physical modeling, the hypothesis that fault systems occupy a gradient of RMT universality classes controlled by geological complexity and fluid involvement (Part~I). We then test this hypothesis against real seismicity from eight global fault zones (Part~II). The test rejects the hypothesis at the mainshock scale---but in doing so reveals a more interesting structure: a \emph{scale-dependent} organization in which the same fault occupies different universality classes at different temporal scales (Part~III).

We present the rejected hypothesis in full rather than discarding it, for two reasons. First, scientific transparency: the forward-modeling predictions define precisely what we expected, which is necessary to appreciate how the real data overturn it. Second, methodological value: the forward model validates the analysis pipeline and quantifies its biases, without which the empirical results could not be properly interpreted. We are explicit throughout that Part~I uses \emph{synthetic} data for hypothesis generation and pipeline validation, not as a source of empirical claims.

The investigation connects to a foundational result: Corral \citep{corral2004} showed that earthquake recurrence-time distributions exhibit universal scaling dominated by clustering, governed by the Omori--Utsu law \citep{utsu1995}. Our contribution is to place this clustering within a unified RMT-classification framework spanning twelve orders of magnitude in time, and to show that the Omori $p$-exponent---once the cluster process is properly isolated---is the non-collapsing physical invariant that the mainshock-repulsion parameter was mistakenly expected to be.

% ═══════════════════════════════════════════════════════════
\part*{Part I --- Hypothesis: A Forward-Modeling Framework}
\addcontentsline{toc}{section}{Part I}
% ═══════════════════════════════════════════════════════════

\section{The Repulsion-Gradient Hypothesis}

We hypothesize three regimes of fault behavior: (i)~geometrically simple, locked strike-slip faults act as stress-cycling oscillators with strong inter-event repulsion ($\beta\to2$); (ii)~fluid-rich subduction interfaces experience pore-pressure modulation that partially decorrelates timing ($\beta\sim1$); (iii)~ancient intraplate fracture zones scatter stress so thoroughly that single-fault memory is erased ($\beta\to0$). The physical reasoning is that the stress shadow following a rupture \citep{stein1999} creates a temporal ``forbidden zone'' analogous to eigenvalue repulsion, which fluids \citep{saffer2003} and structural complexity progressively dismantle.

\section{Forward Modeling and Pipeline Validation}

To test this hypothesis before committing to the data-engineering effort of assembling real catalogs, we built a forward model. For each of eight representative fault zones we generated synthetic catalogs using a Wigner--Dyson point process with a prescribed $\beta_{\rm in}$ assigned from physical reasoning, calibrated to published Gutenberg--Richter rates, then analyzed them blind with the same pipeline later applied to real data. \textbf{These are synthetic catalogs used for hypothesis generation and pipeline validation only.}

\begin{figure}[H]
\centering
\includegraphics[width=0.78\textwidth]{figs/syn_phase_diagram.pdf}
\caption{Synthetic phase diagram (hypothesis). Forward-modeled fault zones span a gradient from GUE-class (simple strike-slip) through GOE (fluid-rich subduction) to the Poisson boundary (complex intraplate). This is the predicted picture to be tested against real data in Part~II.}
\label{fig:synphase}
\end{figure}

Two diagnostics, essential for interpreting the later empirical results, emerge from the forward model. First, the $\langle r\rangle$-based estimator of $\hat\beta$ is substantially biased: it overestimates low $\beta$ and saturates near the GUE ceiling for $\beta\gtrsim1.2$ (Figure~\ref{fig:synrecov}). Consequently only the \emph{ordering} of $\hat\beta$ is reliable, not its absolute value. Second, a power analysis shows that the Kolmogorov--Smirnov (KS) test rejects Poisson with power exceeding 0.99 for sample sizes $n\gtrsim80$; thus, in the empirical analysis, a non-rejection of Poisson at such sample sizes is meaningful rather than a power artifact.

\begin{figure}[H]
\centering
\includegraphics[width=0.6\textwidth]{figs/syn_beta_recovery.pdf}
\caption{Estimator-bias diagnostic from the forward model. Recovered $\hat\beta$ (blue) systematically exceeds the input $\beta$ (dashed line) for low $\beta$ and saturates at high $\beta$. Only the ordering of $\hat\beta$ is trustworthy.}
\label{fig:synrecov}
\end{figure}

\noindent\textbf{Part~I summary.} The forward model demonstrates that, \emph{if} the repulsion-gradient hypothesis were true, our pipeline would detect it (subject to the quantified estimator bias). This establishes a falsifiable prediction---a GOE--Poisson gradient across tectonic regimes---and the diagnostic tools to test it. We now confront this prediction with reality.

% ═══════════════════════════════════════════════════════════
\part*{Part II --- Empirical Test: The Hypothesis Fails at the Mainshock Scale}
\addcontentsline{toc}{section}{Part II}
% ═══════════════════════════════════════════════════════════

\section{Real Data and Processing}

We obtained instrumental catalogs (1970--2024) from the USGS ComCat/ANSS database for the same eight fault zones across three tectonic regimes: strike-slip (Parkfield/San Andreas, North Anatolian Fault, Alpine Fault), subduction (Nankai, Chile, Cascadia), and intraplate (Charlevoix, Longmenshan). Raw catalog sizes ranged from 99 to 3551 events.

Each catalog passed through four stages. (1)~\emph{Magnitude homogenization} to moment magnitude $M_w$ via Scordilis \citep{scordilis2006} global regressions. (2)~\emph{Completeness}: $M_c$ estimated by the maximum-curvature method and applied as a cutoff. (3)~\emph{Declustering} by two independent algorithms---Gardner--Knopoff windowing \citep{gardner1974} and ETAS stochastic declustering \citep{zhuang2002} (200 realizations)---for robustness. (4)~\emph{Spectral statistics}: the normalized spacing distribution, spacing ratio $\langle r\rangle$ \citep{atas2013}, KS tests, and coefficient of variation.

\section{Result: Universal Poisson Collapse}

After declustering, all eight zones yield $\langle r\rangle$ near the Poisson value of 0.386, and in no case is GOE or GUE the preferred fit (Table~\ref{tab:main}, column $\langle r\rangle_{\rm dec}$). Both declustering algorithms agree for every zone. \textbf{The predicted GOE--Poisson gradient is not realized: it collapses to a uniform Poisson plane.} Given the power analysis from Part~I (power $\approx1.0$ at all these sample sizes), this collapse is a genuine rejection of the hypothesis, not a detection failure.

\begin{table}[H]
\centering
\caption{Real-data results across all scales. $\langle r\rangle_{\rm dec}$: declustered mainshock spacing ratio (ETAS). $\langle r\rangle_{\rm full}$, CV, $\gamma$: full-catalog spacing ratio, coefficient of variation, Gamma shape. AF\%: aftershock fraction. $p$: Omori exponent of the isolated cluster process.}
\label{tab:main}
\small
\begin{tabular}{llcccccc}
\toprule
Fault & Regime & $\langle r\rangle_{\rm dec}$ & $\langle r\rangle_{\rm full}$ & CV & $\gamma$ & AF\% & $p$ \\
\midrule
Parkfield & Strike-slip & 0.396 & 0.332 & 2.30 & 0.27 & 62\% & 1.28 \\
NAF & Strike-slip & 0.418 & 0.325 & 1.46 & 0.30 & 48\% & 0.96 \\
Alpine & Strike-slip & 0.454 & 0.374 & 1.90 & 0.38 & 22\% & 0.65 \\
Chile & Subduction & 0.380 & 0.342 & 1.53 & 0.31 & 49\% & 0.88 \\
Nankai & Subduction & 0.426 & 0.370 & 1.08 & 0.55 & 27\% & 0.61 \\
Cascadia & Subduction & 0.368 & 0.338 & 1.58 & 0.31 & 58\% & 0.76 \\
Longmenshan & Intraplate & 0.360 & 0.366 & 3.76 & 0.16 & 79\% & 1.03 \\
Charlevoix & Intraplate & 0.426 & 0.326 & 1.16 & 0.69 & 5\% & -- \\
\bottomrule
\end{tabular}
\end{table}

Two effects explain the collapse. First, declustering by construction removes triggered events, and the residual background of any ETAS-type model is Poisson by definition. Second, each study region aggregates multiple independent fault strands; by the superposition theorem for renewal processes \citep{daley2003}, the sum of many weakly coupled sources tends to Poisson. The elastic-rebound memory of a single segment is diluted by both effects.

\section{A Faint Residual on the Simplest Segment}

Segmenting the NAF along strike, only the eastern segment ($>37^\circ$E)---its simplest, most isolated portion and the initiation point of the 1939 westward rupture sequence---shows a residual repulsion signal ($\hat\beta\approx0.6$--$0.85$ across both declustering methods, $\langle r\rangle\approx0.48$). This signal is not statistically significant ($n=47$; Poisson not rejected), and the classic Parkfield characteristic-earthquake cycle cannot be recovered from the instrumental catalog (only eight $M\geq4.5$ events). Memory, if present at the mainshock scale, survives only as a faint residue on the simplest segments.

% ═══════════════════════════════════════════════════════════
\part*{Part III --- Resolution: Scale-Dependent Universality}
\addcontentsline{toc}{section}{Part III}
% ═══════════════════════════════════════════════════════════

\section{Where the Memory Lives: Clustering, Not Repulsion}

The rejection of the hypothesis is not a dead end. When aftershocks are \emph{retained}, every zone is strongly clustered: the coefficient of variation exceeds unity in all cases (1.08--3.76), and the Gamma shape parameter lies well below unity ($\gamma=0.16$--$0.69$), the signature of a power-law-tailed clustered process (Table~\ref{tab:main}). The full-catalog spacing ratio falls \emph{below} Poisson---the opposite of repulsion (Figure~\ref{fig:acts}). Declustering transforms this clustering into Poisson (Figure~\ref{fig:decl}). The non-trivial temporal structure thus resides in the aftershock cascade, not in mainshock recurrence---consistent with Corral \citep{corral2004}.

\begin{figure}[H]
\centering
\includegraphics[width=\textwidth]{figs/real_three_acts.pdf}
\caption{Left: declustered mainshocks collapse to Poisson ($\langle r\rangle\approx0.39$). Center: full catalogs are clustered (CV$>$1). Right: the Omori $p$-exponent of the isolated cluster process---the non-collapsing invariant.}
\label{fig:acts}
\end{figure}

\begin{figure}[H]
\centering
\includegraphics[width=0.92\textwidth]{figs/real_decluster.pdf}
\caption{Spacing distributions for Parkfield and Cascadia. The full catalog (red) is sharply peaked at $s\to0$ (clustering); declustering (blue) removes the peak, yielding a Poisson-like form.}
\label{fig:decl}
\end{figure}

\section{The Omori Exponent: The Non-Collapsing Invariant}

To probe the cluster process directly, we recorded each aftershock's parent mainshock during Gardner--Knopoff declustering, isolated the aftershock--parent time differences, and fit the modified Omori--Utsu law $n(t)=K/(t+c)^p$ by maximum likelihood \citep{ogata1983}. The resulting decay exponent $p$ varies systematically and does \emph{not} collapse (Table~\ref{tab:main}, last column). The dry, geometrically simple Parkfield segment has the steepest decay ($p=1.28$), whereas fluid-rich subduction interfaces have shallower decays ($p=0.61$--$0.88$), consistent with fluid-driven prolongation of aftershock sequences. Charlevoix produces almost no aftershocks (5\%), too few to fit---itself a meaningful observation about cascade-poor cratonic crust.

This is the key resolution of the investigation. The ``fluid-as-solvent'' intuition of the Part~I hypothesis was correct in spirit but mislocated: fluids do modulate fault dynamics, but their fingerprint is in the aftershock decay rate, not in mainshock recurrence regularity. The physically informative quantity is the Omori $p$-exponent, which does not collapse, rather than a mainshock-repulsion parameter, which does.

\section{Scale-Dependent Universality Classes}

Combining all results with our companion geological study \citep{chen2026geo}, Earth's seismotectonic rhythms occupy three distinct universality classes at three temporal scales (Figure~\ref{fig:scale}):
\begin{itemize}
\item \textbf{Cascade scale (hours--days):} Omori clustering, $\gamma<1$, power-law tails.
\item \textbf{Mainshock scale (years--decades):} Poisson noise, after declustering and spatial superposition.
\item \textbf{Geological-epoch scale (Myr):} GOE repulsion \citep{chen2026geo}, where long-term tectonic reorganizations exhibit weak level repulsion.
\end{itemize}

\begin{figure}[H]
\centering
\includegraphics[width=0.92\textwidth]{figs/scale_schematic.pdf}
\caption{Scale-dependent universality classes spanning twelve orders of magnitude in time. The same Earth system occupies different RMT-related classes at different temporal scales.}
\label{fig:scale}
\end{figure}

\section{Discussion and Conclusions}

This investigation traces a complete scientific arc: a physically motivated hypothesis (Part~I), its rejection by real data (Part~II), and a richer resolution (Part~III). The hypothesis that simple faults exhibit GOE mainshock repulsion is rejected; the genuine, non-collapsing invariant is the Omori $p$-exponent, which varies systematically with tectonic setting and fluid involvement. The scale-dependent picture---clustering at short scales, randomness at intermediate scales, repulsion at geological scales---is richer and more faithful than any single-class description, and unifies our geological and seismological studies within one random-matrix-theoretic framework.

We emphasize the methodological lesson of presenting the rejected hypothesis transparently. The forward model of Part~I was not a source of empirical claims but a tool for generating a falsifiable prediction and validating the analysis pipeline, including the estimator-bias and power diagnostics without which the empirical collapse could be misread. Good science is the arc, not only the endpoint.

\textbf{Limitations.} The Omori fits use aggregate cluster populations rather than individual well-recorded sequences; per-sequence fitting would sharpen the $p$ estimates. The instrumental record is too short to capture characteristic-earthquake cycles for most segments. Magnitude homogenization via global regressions introduces scatter. The residual NAF-east signal warrants targeted study with a denser regional catalog.

\section*{Code and Data Availability}
All code, synthetic and real catalogs, and figures: \url{https://github.com/Ruqing1963/fault-rmt-phase-framework}.\\
Companion geological study (geological-epoch GOE repulsion): \url{https://zenodo.org/records/20766310}.

\begin{thebibliography}{20}
\bibitem[At\'as et al.(2013)]{atas2013} At\'as, Y.Y., Bogomolny, E., Giraud, O., Roux, G. (2013). Distribution of the ratio of consecutive level spacings. \textit{Phys. Rev. Lett.}, 110, 084101.
\bibitem[Chen(2026)]{chen2026geo} Chen, R. (2026). Statistical isomorphism between Phanerozoic tectonic rhythms and Riemann zeta zero spacings. Zenodo, \url{https://zenodo.org/records/20766310}.
\bibitem[Corral(2004)]{corral2004} Corral, \'A. (2004). Long-term clustering, scaling, and universality in the temporal occurrence of earthquakes. \textit{Phys. Rev. Lett.}, 92, 108501.
\bibitem[Daley \& Vere-Jones(2003)]{daley2003} Daley, D.J., Vere-Jones, D. (2003). \textit{An Introduction to the Theory of Point Processes}, Vol. I. Springer.
\bibitem[Gardner \& Knopoff(1974)]{gardner1974} Gardner, J.K., Knopoff, L. (1974). Is the sequence of earthquakes in southern California, with aftershocks removed, Poissonian? \textit{Bull. Seismol. Soc. Am.}, 64, 1363--1367.
\bibitem[Mehta(2004)]{mehta2004} Mehta, M.L. (2004). \textit{Random Matrices}, 3rd ed. Elsevier.
\bibitem[Ogata(1983)]{ogata1983} Ogata, Y. (1983). Estimation of the parameters in the modified Omori formula for aftershock frequencies by the MLE procedure. \textit{J. Phys. Earth}, 31, 115--124.
\bibitem[Reid(1910)]{reid1910} Reid, H.F. (1910). The mechanics of the earthquake. \textit{The California Earthquake of April 18, 1906}, Carnegie Inst.
\bibitem[Saffer \& Bekins(2003)]{saffer2003} Saffer, D.M., Bekins, B.A. (2003). Hydrologic controls on the mechanics of accretionary wedges. \textit{Geology}, 30, 271--274.
\bibitem[Scordilis(2006)]{scordilis2006} Scordilis, E.M. (2006). Empirical global relations converting $M_S$ and $m_b$ to moment magnitude. \textit{J. Seismol.}, 10, 225--236.
\bibitem[Stein(1999)]{stein1999} Stein, R.S. (1999). The role of stress transfer in earthquake occurrence. \textit{Nature}, 402, 605--609.
\bibitem[Utsu et al.(1995)]{utsu1995} Utsu, T., Ogata, Y., Matsu'ura, R.S. (1995). The centenary of the Omori formula. \textit{J. Phys. Earth}, 43, 1--33.
\bibitem[Wigner(1957)]{wigner1957} Wigner, E.P. (1957). Characteristic vectors of bordered matrices. \textit{Ann. Math.}, 62, 548--564.
\bibitem[Zhuang et al.(2002)]{zhuang2002} Zhuang, J., Ogata, Y., Vere-Jones, D. (2002). Stochastic declustering of space-time earthquake occurrences. \textit{J. Am. Stat. Assoc.}, 97, 369--380.
\end{thebibliography}

\end{document}
